Chapter~\ref{chap:distributions} identified three distribution types.
Sequential and near-sequential keys have small local spread, so exact mode
(Section~\ref{sec:adaptive}) applies directly --- FPR~$= 0$ with no hash needed.
When keys are uniform over a large span, truncation
(Section~\ref{sec:truncation}) is effective: large inter-key gaps mean few phantom
collisions and low FPR.

Clustered keys (OSM, geospatial IDs) fit neither case alone: dense regions inside each
cluster behave like near-sequential data, while the sparse gaps between clusters have
large inter-key distances where truncation works well.

The hybrid idea is to split along this structure --- isolate each dense region into its
own exact-mode sub-filter (FPR~$= 0$), and route the sparse keys between clusters to
a fallback filter.

\begin{figure}[h]
    \centering
    \includegraphics[width=0.85\textwidth]{segmentation}
    \caption{Segmented ARE filter: dense clusters go into exact-mode sub-filters
        (FPR~$= 0$); sparse keys between clusters go to a fallback.}
    \label{fig:segmentation}
\end{figure}

This chapter describes Seg-ARE, our segmentation algorithm based on 1D DBSCAN.


\section{1D DBSCAN in $O(n)$}
\label{sec:dbscan-1d}

DBSCAN~\citep{ester1996dbscan} is a density-based clustering algorithm.
It takes two parameters: a neighbourhood width~$\delta$ and a density threshold
\texttt{minPts}. A point is a \emph{core point} if it has at least \texttt{minPts}
neighbours within a window of width~$\delta$; a \emph{border point} if it is within
$\delta$ of a core point but not itself dense; and \emph{noise} otherwise. Clusters
grow from core points by absorbing all density-reachable neighbours.

\begin{figure}[h]
    \centering
    \includegraphics[width=0.85\textwidth]{dbscan_1d}
    \caption{1D DBSCAN on a key sequence (minPts~$= 5$). Blue = core points (window
    contains $\geq$ minPts neighbours). Orange = border (within $\delta$ of a core,
        but $<$ minPts in own window). Gray = noise.}
    \label{fig:dbscan-1d}
\end{figure}

On a finite point set in a general metric space, the naive DBSCAN scan computes the
$\delta$-neighbourhood of each point in $O(n^2)$ time~\citep{ester1996dbscan}; the
original paper recovers $O(n \log n)$ with a spatial index such as an R-tree.
On sorted one-dimensional data, density queries reduce to a sliding window, giving an
$O(n)$ implementation without any auxiliary index.


\section{Seg-ARE: Design and Algorithm}
\label{sec:hybrid-scan}

\paragraph{Border expansion increases bits per key.}
Standard DBSCAN absorbs border points --- points within $\delta$ of a core but not
themselves dense --- into the nearest cluster.
This widens the cluster's span without adding proportional density.
The exact-mode ERE sub-filter encodes keys relative to that universe, so a wider span
means more bits per key.
Seg-ARE drops border expansion: the dense core alone minimises span and keeps the
sub-filter cost low.

\paragraph{Expressing $\delta$ through \texttt{minPts}.}
Standard DBSCAN treats $\delta$ and \texttt{minPts} as independent.
We choose $\delta$ small enough that every range produced by the forward sweep
fits the exact-mode budget of the underlying ERE sub-filter
(see Section~\ref{subsec:fpr-guarantee} for the derivation):

\begin{equation}
    \delta = \texttt{minPts} \cdot \frac{\mathcal{L}}{2\varepsilon} = \frac{256\,\mathcal{L}}{\varepsilon}.
    \label{eq:delta}
\end{equation}

This ties $\delta$ to \texttt{minPts}, leaving a single free parameter.

\paragraph{Choosing \texttt{minPts} = 512.}
The choice balances three concerns:
\begin{itemize}
    \item \textbf{Overhead per key.} Each segment costs roughly 700~bits of fixed metadata
    (\texttt{minKey}, \texttt{maxKey}, pointer, sub-filter header).
    At 512~keys this adds $\sim 1.4$~BPK; for smaller clusters the overhead quickly
    dominates the budget.
    \item \textbf{Query cost.} The segment count is at most $\lfloor n/\texttt{minPts} \rfloor$,
    so the binary search over segment boundaries (Section~\ref{subsec:query}) stays fast
    regardless of dataset size.
    \item \textbf{Detection power.} If \texttt{minPts} is set too high, genuinely dense
    sub-ranges with fewer keys go entirely to the fallback, negating the benefit of
    segmentation.
\end{itemize}
512 balances these concerns: small enough to catch real clusters, large enough to
keep per-segment overhead acceptable.

To sanity-check the scale: for $\mathcal{L} = 128$, $\varepsilon = 10^{-3}$ we get
$\delta \approx 3.3 \times 10^{7}$. Comparing this against the SOSD gap statistics
(Table~\ref{tab:sosd-gaps}), Facebook/Wiki/Books have median gaps well below~$\delta$
so most keys fall into segments, while OSM's median gap of $4.5 \times 10^{7}$ sits
right at the threshold --- segmentation captures only the denser sub-clusters.
A uniform 64-bit distribution has average gap ${\sim}2^{44} \gg \delta$, so every key
goes to the fallback and Seg-ARE degrades to plain truncation.

\subsection{Algorithm}

Seg-ARE runs in three steps:

\begin{enumerate}
    \item \textbf{Forward sweep:} slide a two-pointer window. Whenever the window holds
    $\geq \texttt{minPts}$ keys, record the index range $[\textit{left}, \textit{right}]$.
    \item \textbf{Conditional merge:} adjacent ranges with key-space gap $\leq \delta$ are
    joined only if the combined span is below $m'\mathcal{L}/\varepsilon$,
    where $m'$ is the total key count of the merged segment. Otherwise they remain separate.
    \item \textbf{Fallback:} all keys outside merged ranges go to the fallback filter.
\end{enumerate}

\subsection{Fallback}

Keys not assigned to any segment go to a single fallback ARE filter.
Depending on the density of the fallback keys, this may be a truncation filter
or a hash-based one.

\subsection{Query}\label{subsec:query}

Segments are sorted by $[\texttt{minKey}, \texttt{maxKey}]$. For a query $[a, b]$:
\begin{enumerate}
    \item Binary search over segments intersecting $[a, b]$ --- $O(\log C)$, $C$ = segment
    count.
    \item Check the first intersecting segment only.
    Suppose the query $[a, b]$
    intersects two segments $S_1 = [\texttt{minKey}_1, \texttt{maxKey}_1]$ and
    $S_2 = [\texttt{minKey}_2, \texttt{maxKey}_2]$ with $\texttt{maxKey}_1 < \texttt{minKey}_2$.
    Then $a \leq \texttt{maxKey}_1$ and $b \geq \texttt{minKey}_2$, so the query covers
    $\texttt{maxKey}_1$ --- a real key.
    The exact-mode sub-filter of $S_1$ reports
    non-empty, and no further segments need to be checked.
    \item Check the fallback filter.
    \item Return ``empty'' only if all agree.
\end{enumerate}

\subsection{FPR Guarantee}\label{subsec:fpr-guarantee}

Every segment sub-filter operates in exact mode (FPR~$= 0$).

\textbf{Individual range.} Consider a range of $m$ sorted keys
$k_0 < k_1 < \dots < k_{m-1}$ from the forward sweep.
The window constraint says any \texttt{minPts} consecutive keys span at most
$\delta$: $k_{i+(\texttt{minPts}-1)} - k_i \leq \delta$ for every valid $i$.
We split the $m-1$ gaps into $\lceil(m-1)/(\texttt{minPts}-1)\rceil$
non-overlapping blocks of $\texttt{minPts}-1$ consecutive gaps; each block
spans $\leq \delta$ by the window constraint, so:
\[
    k_{m-1} - k_0 \;\leq\; \left\lceil \frac{m-1}{\texttt{minPts}-1} \right\rceil \cdot \delta.
\]
With $\delta = \texttt{minPts}\cdot\mathcal{L}/(2\varepsilon)$, the ceiling factor
$\lceil(m-1)/(\texttt{minPts}-1)\rceil$ equals~1 for $m \leq \texttt{minPts}$
(span $\leq \delta < m\mathcal{L}/\varepsilon$, trivially satisfied) and first
reaches~2 at $m = \texttt{minPts}+1$, where the margin is tightest:
\[
    2\delta \;=\; \texttt{minPts} \cdot \frac{\mathcal{L}}{\varepsilon}
    \;<\; (\texttt{minPts}+1) \cdot \frac{\mathcal{L}}{\varepsilon}
    \;=\; m \cdot \frac{\mathcal{L}}{\varepsilon}.
\]
Hence every range from the forward sweep fits the exact-mode budget.

\textbf{Merged segments.} The merge step is conditional: two adjacent ranges are
joined only if the combined span is still below $m'\mathcal{L}/\varepsilon$ for the
combined size $m'$. This preserves the exact-mode invariant by construction.

False positives can therefore come only from the fallback filter, giving overall
FPR~$\leq \varepsilon$.

\section{Seg-ARE vs Grafite}
\label{sec:hybrid-ef}
Sections~\ref{sec:hybrid-scan} and~\ref{subsec:fpr-guarantee} have
already told us \emph{when} Seg-ARE is compact: when the input has
dense regions separated by gaps the forward sweep can isolate.
Below we work through two examples and count the bits.

Grafite, when the requested $\varepsilon$ falls below
the lossless threshold $u/n$, degenerates into a single Elias--Fano (EF) encoding of
the raw sorted keys --- the wrapper substitutes
\texttt{ef\_sux\_vector(keys)} and the filter answers range queries by
predecessor probe (Section~\ref{subsec:grafite-fallback}).
Below this threshold Grafite is an approximate filter; above it, it is
essentially EF over the whole universe.
Similar to Adaptive Filter~\ref{sec:adaptive}.

Seg-ARE applies the same EF primitive, but only after segmenting the
input.
Each dense cluster identified by the forward sweep is handed to
an exact-mode ERE sub-filter (Section~\ref{sec:adaptive}), which
encodes the cluster's keys locally; sparse keys outside any cluster
fall through to the truncation hash (Section~\ref{sec:truncation}).
The total size of a Seg-ARE filter is the sum of per-cluster EF
encodings plus the fallback truncation filter --- versus Grafite's
single EF encoding over the whole universe.

This framing makes the win condition concrete: Seg-ARE beats Grafite
exactly when the per-cluster EF plus the fallback together cost less
than one universe-wide EF.
The bpk gap then decomposes as
\[
    \underbrace{\log_2(U_{\text{global}}/n)}_{\text{Grafite EF data}}
    \;-\;
    \underbrace{\Bigl(\sum_{c} \frac{n_c}{n} \cdot
        \log_2\!\bigl(2^{M_c}/n_c\bigr) \;+\; \frac{n_{\text{fb}}}{n} \cdot
        \log_2(1/\varepsilon)\Bigr)}_{\text{Seg-ARE: EF data + truncation fallback}},
\]
where $U_{\text{global}}$ is the spread of the whole input, $2^{M_c}$
and $n_c$ are the local window width and key count of the $c$-th
cluster, and $n_{\text{fb}}$ is the number of keys handed to the
fallback.

The two paragraphs below work through two examples: a synthetic
clustered distribution and the SOSD-Wikipedia timestamps.

\paragraph{Clustered synthetic ($n = 2^{28}$, build-$L = 1$,
$\varepsilon = 10^{-2}$).}
The clustered generator places keys across essentially the full
$2^{64}$ universe, so $\log_2(U/n) = 36$ bits per key for the
universe-wide baseline.

At build-time $\delta = 512 \cdot L / (2\varepsilon) = 25\,600$ the
forward sweep isolates $8911$ dense clusters covering $72.9\%$ of the
keys, with $M_c \in [14.6,\, 28.7]$ bits and
$n_c \in [512,\, 6.7\!\times\!10^{7}]$. The remaining $27.1\%$ fall
through to the truncation hash. The key-weighted per-cluster EF cost
is
\[
    \frac{1}{\sum_c n_c} \sum_c n_c \cdot \log_2\!\Bigl(\frac{2^{M_c}}{n_c}\Bigr)
    \;\approx\; 2.87 \text{ bits per in-cluster key}
\]
--- a $33$-bit reduction in EF \emph{data} cost relative to a
hypothetical universe-wide EF over all $2^{64}$.

Grafite cannot trigger its lossless fallback
(Section~\ref{subsec:grafite-fallback}) at the requested $\varepsilon$
either: lossless requires $n\mathcal{L}/\varepsilon > U$, i.e.\
$\varepsilon < n\mathcal{L}/U = 2^{28}\!\cdot\!128/2^{64} \approx
1.9\!\times\!10^{-9}$, far below the $10^{-2}$ we set. At any
$\varepsilon$ used in this evaluation Grafite stays in approximate
mode and follows the Goswami line $\log_2(\mathcal{L}/\varepsilon) +
O(1)$. Universe-wide lossless EF for this input would actually cost
$\log_2(U/n) + 2 + \text{select-index} \approx 38.94$~BPK
(measured at the saturating end of the Grafite sweep at requested
BPK ${\geq}40$) --- so high that lossless on a $2^{64}$ universe is
useless as a filter: it stores more bits per key than the keys
themselves carry information.

Figure~\ref{fig:seg-clustered-L128} shows the empirical curve at
query length $L = 128$. At matched
$\mathrm{FPR} \approx 6\!\times\!10^{-5}$, Seg-ARE sits at $9.77$ BPK
and Grafite at $20.63$ BPK. Both filters then drop below the
$4\!\times\!10^{-6}$ measurement floor (set by the $10^{6}$-query
batch): Seg-ARE at $10.07$ BPK ($K = 50$) and Grafite at
$26.63$ BPK (requested $\text{BPK}=26$); the true FPR at these points
is non-zero but below detection, not literally zero. The $16.5$-BPK
gap at the floor is therefore a lower bound on the BPK separation:
matching FPRs deeper would widen the gap further, with Grafite's
genuinely zero-FPR point only at the $38.94$~BPK lossless saturation.

The Seg-ARE per-key budget decomposes as
\[
\text{Seg-ARE BPK}\;\approx\;
\underbrace{0.729 \cdot (\,2.87 + 2.53)}_{\text{in-cluster: EF data + ERE metadata}} +
\underbrace{0.271 \cdot w_{\text{fb}}}_{\text{fallback truncation}}
\;=\;3.94 + (\text{fallback}),
\]
with fallback measured (by subtracting the in-cluster contribution
from the observed total) at $10.07 - 3.94 \approx 6.1$~BPK, i.e.\
${\sim}22.6$ bits per fallback key --- consistent with a truncation
filter at $K_{\text{fb}}$ tuned to the same $\varepsilon$ over the
$7.7\!\times\!10^{7}$ leftover keys. Grafite, by contrast, pays
$\log_2(128/6\!\times\!10^{-5}) + {\sim}2 \approx 23$ bits at
$\mathrm{FPR} = 6\!\times\!10^{-5}$, observed at $20.6$ (about $2$
bits below the worst-case constant).

\begin{figure}[h]
    \centering
    \includegraphics[width=0.85\textwidth]{seg_vs_grafite_clustered_L128}
    \caption{FPR vs BPK on the synthetic clustered distribution
    ($n = 2^{28}$, $L = 128$ at query time, $L = 1$ at build time).
    The $16.5$-BPK gap at the $4\!\times\!10^{-6}$ measurement floor
    comes from segmentation. Both filters report empirical
    FPR\,=\,0 at their respective floor points (Seg-ARE at
    $10.07$~BPK, Grafite at $26.63$~BPK); true zero only happens
    once Grafite saturates into universe-wide lossless EF at
    $38.94$~BPK, off-chart on the right.}
    \label{fig:seg-clustered-L128}
\end{figure}

\paragraph{SOSD-Wikipedia ($n = 2^{28}$, $L = 1$,
$\varepsilon = 10^{-2}$).}
Wikipedia edit timestamps form a tight near-sequential cluster:
after deduplication $n \approx 9 \times 10^{7}$ keys span a window
of width $\max - \min \approx 2.37\!\times\!10^{8}$ (about $28$
bits), giving $\log_2(U/n) \approx 1.4$ bits per key for a
universe-wide EF baseline. The segmenter does not collapse the
input into a single cluster: at $\delta = 25\,600$ the forward
sweep finds $359$ dense clusters covering $99.9\%$ of the keys,
with one giant cluster of $7.7\!\times\!10^{7}$ keys
($M_c = 26.5$ bits) and $358$ smaller ones ($M_c$ down to $14.6$
bits). The key-weighted per-cluster EF cost is
\[
    \frac{1}{\sum_c n_c} \sum_c n_c \cdot \log_2\!\Bigl(\frac{2^{M_c}}{n_c}\Bigr)
    \;\approx\; 0.62 \text{ bits per in-cluster key},
\]
$0.77$ bits per key below the universe-wide $1.4$.

Empirically (Figure~\ref{fig:seg-wiki-L128}), Grafite reaches
$\mathrm{FPR} = 0$ at $4.6$ BPK and Seg-ARE at $3.74$ BPK --- a
$0.9$-BPK gap. This matches the segmentation prediction almost
exactly: the $0.77$-bit-per-key drop in EF data cost applied to
the $99.9\%$ in-cluster fraction gives $0.77$ BPK, leaving
${\sim}0.1$ BPK to small implementation constants.

\begin{figure}[h]
    \centering
    \includegraphics[width=0.85\textwidth]{seg_vs_grafite_wiki_L128}
    \caption{FPR vs BPK on SOSD-Wikipedia ($n = 2^{28}$, $L = 128$ at
    query time). Both filters reach $\mathrm{FPR} = 0$ at low
    BPK: Seg-ARE at $3.74$, Grafite at $4.6$. The $0.9$-BPK gap is
    almost entirely the $0.77$ bit-per-key segmentation saving in
    EF data cost (per-cluster $M_c$ values below the universe
    width).}
    \label{fig:seg-wiki-L128}
\end{figure}

The view also predicts the failure modes: when the input has no
isolable dense regions (uniform, OSM), every key falls through to the
fallback, the per-cluster sum collapses to a single truncation-filter
term, and Seg-ARE degrades to plain truncation
(Section~\ref{sec:truncation}). The empirical wins, ties, and losses
across the four SOSD distributions are discussed in
Chapter~\ref{chap:conclusion}.

